Este anexo presenta la fuente de datos, la estructura de los archivos,
los criterios de clasificación sectorial, los niveles de filtrado aplicados
y la caracterización del universo final de acciones utilizado en el modelo P4.
Además, incorpora la metodología actual de análisis con pandemia, en la cual
la pandemia no se elimina como única definición de base, sino que se compara
explícitamente el desempeño de los modelos con y sin las observaciones de
2020--2022 dentro del entrenamiento. \href{https://uccll0.sharepoint.com/:b:/s/uc365_ics2122_g13/IQC6DdeUo4z_Rasgvtcqt8K_AU8x_53YNQQFx14ckIB1xdY?e=DNEyTn}{Ver documento completo}

% ============================================================
\section*{Descripción de la Fuente de Datos}
% ============================================================

Los datos históricos provienen del directorio \texttt{Data/Historical\_Stocks/}
del repositorio del proyecto. Fueron descargados con la librería
\texttt{yfinance} de Yahoo Finance, incluyendo precios de cierre ajustados
por splits, volumen, dividendos y metadatos fundamentales por ticker.

\subsection*{Estructura de archivos}

El directorio contiene \textbf{1.407 archivos}:

\begin{itemize}
  \item \textbf{1.406 archivos CSV} individuales con el formato
        \texttt{stock\_return\_<TICKER>.csv}
  \item \textbf{1 archivo de metadatos} \texttt{stocks\_info.txt} con
        información fundamental de cada ticker (sector, industria,
        descripción, capitalización de mercado, empleados, etc.)
\end{itemize}

\subsection*{Formato de los CSV}

Cada archivo contiene precios diarios con las siguientes columnas:

\begin{center}
\begin{tabular}{lll}
  \toprule
  \textbf{Columna} & \textbf{Tipo} & \textbf{Descripción} \\
  \midrule
  \texttt{Date}         & datetime & Fecha de cierre (zona horaria US/Eastern) \\
  \texttt{Open}         & float    & Precio de apertura ajustado \\
  \texttt{High}         & float    & Precio máximo del día \\
  \texttt{Low}          & float    & Precio mínimo del día \\
  \texttt{Close}        & float    & Precio de cierre ajustado por splits \\
  \texttt{Volume}       & int      & Volumen de transacciones \\
  \texttt{Dividends}    & float    & Dividendo distribuido ese día (USD/acción) \\
  \texttt{Stock Splits} & float    & Factor de split (0 si no hubo split) \\
  \bottomrule
\end{tabular}
\end{center}

\noindent\textbf{Nota:} Los precios están ajustados hacia atrás
(\emph{backward-adjusted}) por splits de acciones, pero \emph{no} por
dividendos. Los dividendos se registran explícitamente en la columna
\texttt{Dividends} y serán usados directamente para alimentar la
\textbf{caja chica} del modelo de optimización.

\subsection*{Cobertura temporal}

\begin{center}
\begin{tabular}{lrr}
  \toprule
  \textbf{Indicador} & \textbf{Valor} & \textbf{Detalle} \\
  \midrule
  Total de tickers en archivos CSV         & 1.406      &              \\
  Tickers con datos válidos                & 1.165      & 83\% del total \\
  Tickers con datos vacíos o errores       & 239        & 17\% sin datos \\
  Fecha de corte más frecuente             & 2026-03-13 & 1.158 tickers  \\
  Mediana de filas por archivo             & 6.532      & $\approx$ 26 años \\
  \midrule
  Tickers con $\geq 5$ años (1.250 filas)  & 1.039      & \\
  Tickers con $\geq 10$ años (2.500 filas) & 912        & \\
  \bottomrule
\end{tabular}
\end{center}

Los años de inicio más frecuentes son 1980 (100 tickers), 1999 (52),
2021 (49) y 1973 (41), lo que refleja una mezcla de empresas consolidadas
y listados recientes, incluyendo SPACs post-2020.

\subsection*{Tratamiento temporal de la pandemia}

La cobertura temporal incluye el periodo de crisis y recuperación asociado a
COVID-19. En versiones preliminares del análisis, la pandemia se trataba
principalmente como un periodo a excluir. En la versión actual, el criterio se
modificó: se construyen dos escenarios comparables, uno que excluye 2020--2022
de la ventana de entrenamiento y otro que conserva ese periodo.

Esta decisión permite medir si los datos de estrés alteran las estimaciones de
retorno, volatilidad y covarianza, sin confundir ese efecto con un cambio en el
periodo futuro evaluado. Por lo tanto, la pandemia pasa de ser solo un filtro
de limpieza a ser una variable metodológica explícita del diseño experimental.

% ============================================================
\section*{Clasificación Sectorial}
% ============================================================

\subsection*{Sistema de clasificación}

Los tickers se clasifican usando la taxonomía sectorial de
\textbf{GICS}\footnote{Global Industry Classification Standard,
desarrollado por MSCI y S\&P Dow Jones Indices.}
a través de los metadatos de Yahoo Finance, que asigna a cada empresa
un \texttt{sector} y una \texttt{industry} (subsector).

El archivo \texttt{stocks\_info.txt} contiene 1.406 entradas en formato:

\begin{center}
  \texttt{TICKER;\{`sector': `...', `industry': `...', `marketCap': ..., ...\}}
\end{center}

\subsection*{Distribución sectorial del universo completo}

La Tabla~\ref{tab:sectores} muestra la distribución de los tickers con datos
válidos por sector.

\begin{table}[H]
  \centering
  \caption{Distribución sectorial del universo completo con datos válidos}
  \label{tab:sectores}
  \begin{tabular}{lrr}
    \toprule
    \textbf{Sector} & \textbf{Tickers} & \textbf{\% del total} \\
    \midrule
    Financial Services      & 716 & 61.4\% \\
    Technology              &  81 &  6.9\% \\
    Industrials             &  74 &  6.4\% \\
    Healthcare              &  64 &  5.5\% \\
    Consumer Cyclical       &  59 &  5.1\% \\
    Consumer Defensive      &  38 &  3.3\% \\
    Real Estate             &  32 &  2.7\% \\
    Utilities               &  31 &  2.7\% \\
    Basic Materials         &  22 &  1.9\% \\
    Energy                  &  21 &  1.8\% \\
    Communication Services  &  20 &  1.7\% \\
    Unknown                 &   8 &  0.7\% \\
    \midrule
    \textbf{Total}          & \textbf{1.166} & \textbf{100\%} \\
    \bottomrule
  \end{tabular}
\end{table}

\noindent\textbf{Nota:} Financial Services representa el 61\% del universo crudo
debido a la alta presencia de bancos regionales, gestores de activos y
shell companies (SPACs). El filtrado posterior reducirá significativamente
este sesgo sectorial.

\subsection*{Subsectores predominantes en Financial Services}

\begin{center}
\begin{tabular}{lrr}
  \toprule
  \textbf{Subsector (Industry)} & \textbf{Tickers} & \textbf{\%} \\
  \midrule
  Banks -- Regional                  & 317 & 44.3\% \\
  Asset Management                   & 101 & 14.1\% \\
  Capital Markets                    &  56 &  7.8\% \\
  Credit Services                    &  45 &  6.3\% \\
  Shell Companies                    &  42 &  5.9\% \\
  Insurance -- Property \& Casualty  &  41 &  5.7\% \\
  Banks -- Diversified               &  19 &  2.7\% \\
  Insurance -- Life                  &  17 &  2.4\% \\
  Insurance -- Specialty             &  16 &  2.2\% \\
  Insurance Brokers                  &  16 &  2.2\% \\
  Otros                              &  46 &  6.4\% \\
  \midrule
  \textbf{Total}                     & \textbf{716} & \textbf{100\%} \\
  \bottomrule
\end{tabular}
\end{center}

\subsection*{Tipos de instrumento}

El campo \texttt{quoteType} de Yahoo Finance indica que 1.165 tickers
corresponden a \texttt{EQUITY} (acciones individuales de empresas) y
1 ticker es \texttt{MUTUALFUND} (excluido del análisis). El universo
es prácticamente homogéneo en acciones individuales.

% ============================================================
\section*{Capitalización de Mercado}
% ============================================================

Se definen los siguientes rangos estándar de capitalización:

\begin{center}
\begin{tabular}{llrr}
  \toprule
  \textbf{Categoría} & \textbf{Rango (USD)} & \textbf{Tickers} & \textbf{\%} \\
  \midrule
  Mega cap   & $\geq$ 200B    &  50 &  4.3\% \\
  Large cap  & 10B -- 200B    & 473 & 40.6\% \\
  Mid cap    & 2B -- 10B      & 193 & 16.6\% \\
  Small cap  & 300M -- 2B     & 234 & 20.1\% \\
  Micro cap  & $<$ 300M       & 210 & 18.0\% \\
  \midrule
  Sin dato   & ---            &   6 &  0.5\% \\
  \midrule
  \textbf{Total con datos} & & \textbf{1.160} & \\
  \bottomrule
\end{tabular}
\end{center}

La capitalización de mercado se utiliza en dos niveles. Primero, permite
filtrar empresas demasiado pequeñas o sin información suficiente, que podrían
introducir problemas de liquidez o sesgos extremos en la estimación. Segundo,
sirve como criterio para construir referencias de mercado, en particular el
prior de equilibrio de Black-Litterman.

% ============================================================
\section*{Niveles de Filtrado}
% ============================================================

Se definen cinco niveles de filtrado en cascada, cada uno más restrictivo
que el anterior. El objetivo es obtener un universo operacional que sea
\textbf{líquido, estable, representativo y estadísticamente válido}
para el optimizador de Markowitz.

La secuencia de filtros es la siguiente:

\begin{itemize}
  \item \textbf{F0 --- Universo crudo:} 1.406 tickers
  \item \textbf{F1 --- Liquidez básica:} 883 tickers (precio $\geq$ 5 USD, historia $\geq$ 5 años)
  \item \textbf{F2 --- Sin shells ni desconocidos:} 881 tickers (excluye Shell Companies y sector Unknown)
  \item \textbf{F3 --- Large/Mega cap:} 686 tickers (capitalización $\geq$ 2B USD)
  \item \textbf{F4 --- Historia suficiente:} 641 tickers ($\geq$ 10 años de datos)
  \item \textbf{F5 --- Volatilidad razonable (universo final descriptivo):} 636 tickers ($5\% \leq \sigma_{ann} \leq 100\%$)
\end{itemize}

\subsection*{Razones de exclusión}

La Tabla~\ref{tab:exclusion} desglosa cuántos tickers fueron eliminados
en cada filtro y por qué razón:

\begin{table}[H]
  \centering
  \caption{Razones de exclusión del universo final F5}
  \label{tab:exclusion}
  \begin{tabular}{lrr}
    \toprule
    \textbf{Razón de exclusión} & \textbf{Tickers} & \textbf{\%} \\
    \midrule
    Historial $< 10$ años             & 254 & 47.9\% \\
    Cap.\ mkt $< 2$B o sin dato       & 233 & 43.9\% \\
    Precio de cierre $< 5$ USD        &  37 &  7.0\% \\
    Volatilidad fuera de rango        &   5 &  0.9\% \\
    Datos vacíos o errores de lectura & 239 & ---    \\
    \midrule
    \textbf{Total excluidos (F0$\to$F5)} & \textbf{529} & \\
    \bottomrule
  \end{tabular}
\end{table}

\subsection*{Definición formal de cada filtro}

\textbf{F1 --- Liquidez básica.}
Un ticker $i$ pertenece a $\mathcal{F}_1$ si y solo si cumple:
número de filas con precios válidos $n_i \geq 1.250$ ($\approx$ 5 años de historia),
precio de cierre más reciente $P_i^{\text{last}} \geq 5$ USD, y
capitalización de mercado $\text{Cap}_i \geq 300$M USD (o sin dato).
Esto excluye penny stocks y acciones ilíquidas sin historia suficiente
para estimar parámetros estadísticos confiables.

\medskip

\textbf{F2 --- Calidad de clasificación.}
$\mathcal{F}_2 \subset \mathcal{F}_1$: adicionalmente se excluye si
\texttt{industry} = ``Shell Companies'' o \texttt{sector} = ``Unknown''.
Las shell companies, principalmente SPACs, no tienen operaciones reales
ni retornos representativos, y los tickers sin sector asignado no pueden
incluirse en análisis sectoriales.

\medskip

\textbf{F3 --- Large/Mega cap.}
$\mathcal{F}_3 \subset \mathcal{F}_2$: adicionalmente se requiere
$\text{Cap}_i \geq 2 \times 10^9$ USD.
Las empresas de capitalización media-alta tienen mayor liquidez de mercado,
menores spreads bid-ask, y sus precios reflejan mejor la información pública
disponible, condición implícita del modelo de Markowitz.

\medskip

\textbf{F4 --- Historia suficiente.}
$\mathcal{F}_4 \subset \mathcal{F}_3$: adicionalmente $n_i \geq 2.500$
($\approx$ 10 años de datos diarios).
Con 10 años de historia se captura al menos un ciclo económico completo,
incluyendo crisis (2008--09, 2020) y expansiones, lo que es necesario
para la estimación robusta de matrices de covarianza.

\medskip

\textbf{F5 --- Volatilidad razonable (universo final descriptivo).}
$\mathcal{F}_5 \subset \mathcal{F}_4$: adicionalmente se requiere:
\[
  5\% \;\leq\; \hat{\sigma}_{i}^{\text{ann}} \;\leq\; 100\%
  \qquad \text{donde} \qquad
  \hat{\sigma}_{i}^{\text{ann}} = \hat{\sigma}_{i}^{\text{diaria}} \times \sqrt{252}
\]
$\sigma < 5\%$ indicaría datos incorrectos o instrumentos no bursátiles;
$\sigma > 100\%$ corresponde a activos especulativos extremos incompatibles
con el modelo de media-varianza.

% ============================================================
\section*{Universo Final Descriptivo: 636 Acciones}
% ============================================================

El universo final descriptivo contiene \textbf{636 acciones} con historia
$\geq$ 10 años, capitalización $\geq$ 2B USD, precio $\geq$ 5 USD, y
volatilidad anualizada entre 5\% y 100\%.

Se denomina ``descriptivo'' porque resume el resultado completo del proceso de
depuración sobre los datos disponibles. Para los experimentos con y sin
pandemia se aplica una restricción adicional de comparabilidad temporal,
explicada más adelante, que reduce el universo operacional a 601 acciones.

\subsection*{Distribución sectorial del universo final}

\begin{table}[H]
  \centering
  \caption{Estadísticas por sector --- universo F5 (636 tickers)}
  \label{tab:sectorstats}
  \begin{tabular}{lrrrr}
    \toprule
    \textbf{Sector} & \textbf{N} & \textbf{$\bar{\sigma}_{ann}$}
                    & \textbf{$\overline{\text{CAGR}}$} & \textbf{\% con dividendos} \\
    \midrule
    Financial Services     & 237 & 37.5\% & 10.6\% & 95\% \\
    Technology             &  73 & 44.6\% & 16.5\% & 70\% \\
    Industrials            &  64 & 34.6\% & 13.2\% & 95\% \\
    Healthcare             &  59 & 37.3\% & 13.4\% & 66\% \\
    Consumer Cyclical      &  55 & 41.2\% & 13.9\% & 78\% \\
    Consumer Defensive     &  33 & 29.3\% & 11.4\% & 97\% \\
    Utilities              &  29 & 24.7\% &  9.7\% & 100\% \\
    Real Estate            &  29 & 34.8\% & 11.3\% & 93\% \\
    Energy                 &  20 & 37.5\% & 11.5\% & 100\% \\
    Basic Materials        &  19 & 35.1\% & 11.5\% & 100\% \\
    Communication Services &  18 & 36.7\% & 14.6\% & 78\% \\
    \midrule
    \textbf{Total}         & \textbf{636} & \textbf{36.6\%} & \textbf{11.9\%} & \textbf{90\%} \\
    \bottomrule
  \end{tabular}
\end{table}

Utilities es el sector de menor volatilidad promedio (24.7\%), adecuado
para perfiles conservadores con dividendos garantizados (100\%).
Technology lidera en CAGR (16.5\%) pero con alta volatilidad (44.6\%).
Consumer Defensive combina volatilidad baja (29.3\%) con dividendos casi
universales (97\%). Financial Services sigue siendo el sector más grande
(37\% del universo F5).

\begin{table}[H]
  \centering
  \caption{Top subsectores por sector en el universo F5}
  \label{tab:subsectores}
  \scriptsize
  \setlength{\tabcolsep}{2pt}
  \renewcommand{\arraystretch}{0.82}
  \begin{adjustbox}{max width=0.78\textwidth}
  \begin{tabular}{@{}llr@{}}
    \toprule
    \textbf{Sector} & \textbf{Subsector (Industry)} & \textbf{N} \\
    \midrule
    \multirow{5}{*}{Financial Services}
      & Banks -- Regional               & 93 \\
      & Asset Management                & 26 \\
      & Capital Markets                 & 19 \\
      & Insurance -- P\&C               & 18 \\
      & Banks -- Diversified            & 18 \\
    \midrule
    \multirow{5}{*}{Technology}
      & Semiconductors                  & 14 \\
      & Software -- Application         & 13 \\
      & Software -- Infrastructure      & 12 \\
      & IT Services                     & 10 \\
      & Computer Hardware               &  5 \\
    \midrule
    \multirow{4}{*}{Industrials}
      & Specialty Industrial Machinery  & 14 \\
      & Aerospace \& Defense            & 11 \\
      & Integrated Freight \& Logistics &  5 \\
      & Airlines                        &  4 \\
    \midrule
    \multirow{4}{*}{Healthcare}
      & Diagnostics \& Research         & 11 \\
      & Medical Devices                 & 10 \\
      & Drug Manufacturers -- General   &  9 \\
      & Medical Instruments \& Supplies &  9 \\
    \midrule
    \multirow{4}{*}{Consumer Cyclical}
      & Packaging \& Containers         &  6 \\
      & Auto Parts                      &  6 \\
      & Restaurants                     &  6 \\
      & Travel Services                 &  5 \\
    \midrule
    \multirow{3}{*}{Real Estate}
      & REIT -- Specialty               &  7 \\
      & REIT -- Residential             &  6 \\
      & REIT -- Retail                  &  5 \\
    \midrule
    \multirow{3}{*}{Energy}
      & Oil \& Gas E\&P                 &  8 \\
      & Oil \& Gas Midstream            &  4 \\
      & Oil \& Gas Refining             &  3 \\
    \midrule
    \multirow{2}{*}{Basic Materials}
      & Specialty Chemicals             & 10 \\
      & Agricultural Inputs             &  2 \\
    \midrule
    \multirow{3}{*}{Utilities}
      & Utilities -- Regulated Electric & 23 \\
      & Utilities -- Diversified        &  2 \\
      & Utilities -- Regulated Gas      &  2 \\
    \midrule
    \multirow{3}{*}{Communication Services}
      & Entertainment                   &  6 \\
      & Telecom Services                &  5 \\
      & Internet Content \& Information &  4 \\
    \bottomrule
  \end{tabular}
  \end{adjustbox}
\end{table}

\subsection*{Representantes por sector (top 5 por capitalización)}

\begin{table}[H]
  \centering
  \caption{Top 5 Technology (por capitalización de mercado)}
  \small
  \begin{tabular}{llrrrr}
    \toprule
    \textbf{Ticker} & \textbf{Empresa} & \textbf{Cap (USD B)}
                    & \textbf{$\sigma_{ann}$} & \textbf{CAGR} & \textbf{Inicio} \\
    \midrule
    NVDA & NVIDIA Corp.      & 4.381 & 59.5\% & 36.7\% & 1999 \\
    AAPL & Apple Inc.        & 3.676 & 43.8\% & 18.9\% & 1980 \\
    MSFT & Microsoft Corp.   & 2.940 & 33.2\% & 24.6\% & 1986 \\
    AVGO & Broadcom Inc.     & 1.527 & 37.9\% & 40.6\% & 2009 \\
    MU   & Micron Technology &   480 & 60.3\% & 14.8\% & 1984 \\
    \bottomrule
  \end{tabular}
\end{table}

\begin{table}[H]
  \centering
  \caption{Top 5 Healthcare (por capitalización de mercado)}
  \small
  \begin{tabular}{llrrrr}
    \toprule
    \textbf{Ticker} & \textbf{Empresa} & \textbf{Cap (USD B)}
                    & \textbf{$\sigma_{ann}$} & \textbf{CAGR} & \textbf{Inicio} \\
    \midrule
    LLY  & Eli Lilly \& Co.   & 882 & 27.5\% & 13.9\% & 1972 \\
    JNJ  & Johnson \& Johnson & 582 & 22.8\% & 13.8\% & 1962 \\
    ABBV & AbbVie Inc.        & 388 & 26.3\% & 19.6\% & 2013 \\
    MRK  & Merck \& Co.       & 286 & 24.9\% & 12.4\% & 1962 \\
    UNH  & UnitedHealth Group & 256 & 39.8\% & 20.9\% & 1984 \\
    \bottomrule
  \end{tabular}
\end{table}

\begin{table}[H]
  \centering
  \caption{Top 5 Financial Services (por capitalización de mercado)}
  \small
  \begin{tabular}{llrrrr}
    \toprule
    \textbf{Ticker} & \textbf{Empresa} & \textbf{Cap (USD B)}
                    & \textbf{$\sigma_{ann}$} & \textbf{CAGR} & \textbf{Inicio} \\
    \midrule
    JPM  & JPMorgan Chase  & 764 & 35.0\% & 13.0\% & 1980 \\
    V    & Visa Inc.       & 592 & 28.5\% & 19.5\% & 2008 \\
    MA   & Mastercard Inc. & 444 & 32.4\% & 27.4\% & 2006 \\
    BAC  & Bank of America & 341 & 37.5\% &  6.7\% & 1973 \\
    HSBC & HSBC Holdings   & 269 & 27.0\% &  5.9\% & 1999 \\
    \bottomrule
  \end{tabular}
\end{table}

\begin{table}[H]
  \centering
  \caption{Top 5 Consumer Defensive + Utilities (perfiles conservadores)}
  \small
  \begin{tabular}{lllrrrr}
    \toprule
    \textbf{Ticker} & \textbf{Sector} & \textbf{Empresa} & \textbf{Cap (B)}
                    & \textbf{$\sigma_{ann}$} & \textbf{CAGR} & \textbf{Inicio} \\
    \midrule
    WMT  & Cons.\ Def. & Walmart Inc.      & 1.009 & 29.1\% & 19.0\% & 1972 \\
    COST & Cons.\ Def. & Costco Wholesale  &   447 & 31.0\% & 13.2\% & 1986 \\
    PG   & Cons.\ Def. & Procter \& Gamble &   353 & 21.3\% & 10.5\% & 1962 \\
    KO   & Cons.\ Def. & Coca-Cola Co.     &   333 & 22.8\% & 12.3\% & 1962 \\
    NEE  & Utilities   & NextEra Energy    &   193 & 21.8\% & 11.1\% & 1973 \\
    \bottomrule
  \end{tabular}
\end{table}

% ============================================================
\section*{Universo Operacional con Inclusión de Pandemia: 601 Acciones}
% ============================================================

Para los experimentos de portafolio no se usa directamente la totalidad del
universo descriptivo de 636 acciones. Se requiere que los activos sean
comparables entre los escenarios \texttt{sin\_pandemia} y \texttt{con\_pandemia}
y que tengan datos suficientes en las ventanas de entrenamiento y evaluación.
Al aplicar esta intersección temporal y operacional, el universo usado por los
notebooks principales queda en \textbf{601 acciones}.

Esta diferencia no contradice el universo F5. El número 636 describe el
resultado completo del proceso de filtrado de datos; el número 601 corresponde
al subconjunto común y operacional usado para comparar metodologías bajo el
horizonte \texttt{h7\_f2}.

\subsection*{Escenarios de datos}

La metodología actual define dos escenarios:

\begin{itemize}
  \item \textbf{\texttt{sin\_pandemia}:} excluye las observaciones de 2020--2022
        de la base usada para calibrar los modelos.
  \item \textbf{\texttt{con\_pandemia}:} conserva las observaciones de 2020--2022
        dentro del historial de entrenamiento.
\end{itemize}

En ambos escenarios se mantiene fijo el periodo futuro fuera de muestra:
2024-01-29 a 2026-01-30. Por lo tanto, el experimento no compara dos futuros
distintos. Compara dos formas de estimar los parámetros con distinta
información histórica.

\begin{table}[H]
  \centering
  \caption{Diagnóstico de escenarios para el universo operacional}
  \label{tab:anexo2_escenarios_pandemia}
  \small
  \begin{tabular}{lrrrr}
    \toprule
    \textbf{Escenario} & \textbf{Activos} & \textbf{Retornos comunes}
    & \textbf{Días pandemia train} & \textbf{Días pandemia test} \\
    \midrule
    Sin pandemia & 601 & 2.489 & 0 & 0 \\
    Con pandemia & 601 & 3.245 & 756 & 0 \\
    \bottomrule
  \end{tabular}
\end{table}

\begin{table}[H]
  \centering
  \caption{Ventanas temporales en el experimento h7\_f2}
  \label{tab:anexo2_ventanas_pandemia}
  \small
  \begin{adjustbox}{max width=\textwidth}
  \begin{tabular}{llll}
    \toprule
    \textbf{Escenario} & \textbf{Retornos disponibles}
    & \textbf{Entrenamiento} & \textbf{Evaluación futura} \\
    \midrule
    Sin pandemia & 2013-03-08 a 2026-01-30
      & 2014-01-23 a 2024-01-26
      & 2024-01-29 a 2026-01-30 \\
    Con pandemia & 2013-03-08 a 2026-01-30
      & 2017-01-24 a 2024-01-26
      & 2024-01-29 a 2026-01-30 \\
    \bottomrule
  \end{tabular}
  \end{adjustbox}
\end{table}

La ventana de entrenamiento del escenario sin pandemia comienza antes porque
necesita recorrer más tiempo calendario para reunir una cantidad comparable de
observaciones tras remover 2020--2022. En cambio, el escenario con pandemia
conserva esos días y puede completar la muestra de entrenamiento con una
ventana calendario más corta.

\subsection*{Interpretación metodológica}

El tratamiento con/sin pandemia permite separar tres efectos:

\begin{itemize}
  \item \textbf{Efecto de limpieza:} se mantiene un universo líquido,
        clasificado y con historia suficiente.
  \item \textbf{Efecto de calibración:} se observa cómo cambian $\mu$,
        $\Sigma$, la frontera eficiente y los pesos cuando se incorpora un
        periodo de estrés.
  \item \textbf{Efecto de evaluación:} se mantiene fijo el futuro fuera de
        muestra, por lo que las diferencias no se atribuyen a distintos
        mercados futuros.
\end{itemize}

Así, la pandemia no se interpreta como ruido que deba eliminarse
automáticamente. Se interpreta como información histórica extrema cuya
inclusión puede mejorar o deteriorar la estimación, dependiendo del modelo.

\subsection*{Efecto observado en resultados}

La inclusión de 2020--2022 afecta principalmente a los portafolios que dependen
de retornos esperados históricos. La Tabla~\ref{tab:anexo2_resultados_pandemia}
resume el efecto en el experimento Markowitz vs Benchmark.

\begin{table}[H]
  \centering
  \caption{Efecto de incluir pandemia en el entrenamiento}
  \label{tab:anexo2_resultados_pandemia}
  \small
  \begin{adjustbox}{max width=\textwidth}
  \begin{tabular}{lrrrrr}
    \toprule
    \textbf{Portafolio} & \textbf{Ret. sin} & \textbf{Sharpe sin}
    & \textbf{Ret. con} & \textbf{Sharpe con} & \textbf{$\Delta$ Sharpe} \\
    \midrule
    Benchmark & 101,4\% & 2,05 & 101,4\% & 2,05 & 0,00 \\
    Mínima varianza & 56,6\% & 2,11 & 54,6\% & 2,02 & $-$0,09 \\
    Muy conservador & 58,0\% & 2,21 & 54,5\% & 2,02 & $-$0,18 \\
    Conservador & 67,8\% & 2,42 & 55,5\% & 1,97 & $-$0,45 \\
    Neutro & 82,9\% & 2,41 & 60,5\% & 1,76 & $-$0,66 \\
    Ideal de mercado & 76,5\% & 2,51 & 59,5\% & 1,40 & $-$1,11 \\
    Arriesgado & 91,6\% & 1,73 & 58,4\% & 1,17 & $-$0,57 \\
    Muy arriesgado & 102,7\% & 1,41 & 44,0\% & 0,55 & $-$0,86 \\
    \bottomrule
  \end{tabular}
  \end{adjustbox}
\end{table}

El benchmark permanece invariante porque su regla no depende de la estimación
de $\mu$ ni de $\Sigma$. En cambio, las carteras optimizadas sí se modifican.
El deterioro más fuerte ocurre en Ideal de mercado ($-1,11$ puntos de Sharpe)
y Muy arriesgado ($-0,86$), lo que muestra que los perfiles más agresivos son
los más sensibles a la inclusión de datos de crisis. Mínima varianza cae solo
$-0,09$ puntos de Sharpe, porque depende más de covarianzas que de retornos
esperados individuales.

% ============================================================
\section*{Implicaciones para el Modelo de Optimización}
% ============================================================

\subsection*{Conexión con parámetros del modelo}

El universo F5 alimenta directamente los parámetros del modelo
MO-MS-SNLP descrito en la formulación matemática del P4:

\begin{table}[H]
  \centering
  \caption{Conexión entre parámetros del modelo y fuente de datos}
  \label{tab:parametros_modelo}
  \footnotesize
  \setlength{\tabcolsep}{3pt}
  \renewcommand{\arraystretch}{0.95}
  \begin{tabularx}{0.92\textwidth}{@{}l>{\raggedright\arraybackslash}X>{\raggedright\arraybackslash}X@{}}
    \toprule
    \textbf{Parámetro} & \textbf{Fuente en los datos} & \textbf{Notas} \\
    \midrule
    $r_{i,t}^s$ (retorno escenario $s$)
      & Columna \texttt{Close}, retornos diarios
      & Agregado a retornos semanales \\
    $d_{i,t}$ (dividendos)
      & Columna \texttt{Dividends}
      & Alimenta \texttt{CC\_t} (caja chica) \\
    $\mu_i$ (retorno esperado)
      & CAGR histórico o media móvil
      & Percentil 10/50/90 por escenario \\
    $\Sigma_{ij}$ (covarianza)
      & Retornos diarios, ventana de 2 años
      & Estimación rolling para rebalanceo \\
    $\alpha_p$ (tolerancia de pérdida)
      & Perfil de riesgo del cliente
      & 5 niveles: 0\%, 5\%, 15\%, 30\%, 40\% \\
    \bottomrule
  \end{tabularx}
\end{table}

En la implementación actual, estos parámetros se calculan sobre dos bases
comparables: una sin pandemia y otra con pandemia. Por lo tanto, el mismo
pipeline de datos permite medir cómo la inclusión de 2020--2022 cambia los
insumos de optimización.

\subsection*{Sub-universos recomendados por perfil de riesgo}

Para mantener el problema computacionalmente tratable y respetar
la restricción de respuesta $< 10$ segundos, se sugieren los siguientes
sub-universos:

\begin{table}[H]
  \centering
  \caption{Sub-universos recomendados por perfil de riesgo}
  \label{tab:subuniverse}
  \begin{tabular}{llp{6.5cm}}
    \toprule
    \textbf{Perfil} & \textbf{Tamaño sugerido} & \textbf{Criterio de selección} \\
    \midrule
    Muy conservador (0\%)  & 30--50 acciones   & $\sigma_{ann} < 30\%$, dividendos, sectores Utilities y Cons.\ Def. \\
    Conservador (5\%)      & 50--80 acciones   & $\sigma_{ann} < 35\%$, dividendos $\geq 50\%$ del universo sectorial \\
    Neutro (15\%)          & 80--120 acciones  & Todos los sectores, pesos del universo F5 \\
    Arriesgado (30\%)      & 100--150 acciones & Incluye Tech y Consumer Cyclical de alta volatilidad \\
    Muy arriesgado (40\%)  & Sin restricción   & Universo F5 completo en descriptivo; universo operacional de 601 en experimentos con/sin pandemia \\
    \bottomrule
  \end{tabular}
\end{table}

Para el prototipo inicial se recomienda trabajar con un universo de
\textbf{50--100 acciones} seleccionadas por sector (representantes por
capitalización), lo que garantiza la resolución en Gurobi en $< 1$
segundo para el LP de capas 0--2.

\subsection*{Consideraciones sobre la caja chica}

Los dividendos en la columna \texttt{Dividends} representan el monto
en USD por acción en la fecha de pago. Para el modelo:

\begin{align}
  d_{i,t} &= \text{Dividends}_{i,t} \cdot \frac{w_{i,t} \cdot V_{t-1}}{P_{i,t}}
\end{align}

donde $P_{i,t}$ es el precio de cierre de la semana $t$ y $V_{t-1}$
el valor del portafolio en la semana anterior. El 90\% del universo F5
distribuye dividendos, lo que implica que la caja chica será activa
para la gran mayoría de portafolios.

\subsection*{Implicancia de la pandemia para los parámetros}

La pandemia afecta principalmente la estimación de retornos esperados y
covarianzas. En Markowitz, esto se traduce en cambios directos en la frontera
eficiente y en la selección de portafolios por perfil. En Black-Litterman, el
prior de equilibrio y las \textit{views} amortiguan parcialmente la dependencia
de la media histórica, pero la matriz de covarianzas y la volatilidad siguen
capturando parte del shock. En Simulación Monte Carlo, los KPIs resultantes se
usan como insumo para proyectar riqueza, retiro del cliente y utilidad de la
empresa.

Por ello, mantener ambos escenarios es metodológicamente más informativo que
reportar una sola base sin pandemia. Permite observar si una técnica es robusta
cuando se entrena con datos de estrés y si el resultado depende de un supuesto
de limpieza demasiado favorable.

% ============================================================
\section*{Conclusión}
% ============================================================

En síntesis, a partir de un universo inicial de \textbf{1.406 tickers}
se construyó un universo final descriptivo de \textbf{636 acciones} aptas para
alimentar el modelo de optimización del proyecto. El proceso de depuración
permitió eliminar instrumentos con información incompleta, baja liquidez,
escasa historia o comportamiento incompatible con el enfoque de media-varianza.

El resultado es un universo más estable, representativo y utilizable
computacionalmente, con presencia de 11 sectores económicos, una volatilidad
promedio anualizada de 36.6\%, y una alta proporción de acciones que pagan
dividendos. Esto permite que el conjunto final no solo sea adecuado para la
estimación de retornos y covarianzas, sino también consistente con la lógica
financiera del modelo, incluyendo la incorporación de caja chica vía dividendos.

La metodología actual agrega una capa adicional: para los experimentos de
Markowitz, Black-Litterman y Simulación Monte Carlo se trabaja con un universo
operacional comparable de \textbf{601 acciones}, evaluado bajo escenarios
\texttt{sin\_pandemia} y \texttt{con\_pandemia}. En ambos casos el periodo
futuro se mantiene fijo, por lo que la comparación mide el efecto de incluir
2020--2022 en la calibración, no el efecto de cambiar el mercado evaluado.

\vspace{1em}
\begin{center}
  \small
  Documento generado con datos al 13 de marzo de 2026.\\
  Fuente: Yahoo Finance vía \texttt{yfinance}. Análisis: Python 3.11.
\end{center}
